Opgaven, De skal læse, omhandler spillet \textit{Project Plastic Waste}. Det har ikke hele tiden heddet sådan, så hvis det undervejs ikke nævnes ved navn skyldes det dette. Rapporten er opdelt på en efter, hvad der er vurderet mest hensigtmæssigt for en læser. Kapitler forekommer i deres naturlige progression; således kan De forvente, at den spæde opstart, spillets vision og overvejelser kommer først i rapporten, dernæst selve produktudviklingen, og afrundingsvis en evaluering af produktet, læren ved processen, en række forslag til videreudvikling og afrundende bemærkninger. Allerbagerst i opgaven findes biliografien og dernæst bilagene. For en forståelse af projektets grundlag henvises der til projektoplægget (\autoref{sec:projektoplæg}), men i allergroveste form er der tale om et spil, der har til formål, at formidle et økologisk budskab, hvori opsamling af skrald spiller en essentiel rolle. Om økologisme, kan De læse om i problemanalysen (\autoref{sec:problemanalyse}). For at kunne forstå, hvad spillet grundlæggende går ud på, henvises der til spillets elevator-pitch (\autoref{par:elevator_pitch}), der også kommer snarligt herefter. Det kan være en fordel at se den vedhæftede playthrough-trailer med det samme, og der opfordres til at man bidder særligt mærke i afsnittet om Bevy-spilmotoren og det datadrevne ECS-paradigme (\autoref{par:bevy}), som motoren bygger på, da særlig terminologi kort drøftes og denne muligvis er ukendt for læseren.

Ellers kan rapporten sagtens læses i kronologisk rækkefølge. God læsning.
